\item \textbf{Modos longitudinales en cavidad láser y efecto Doppler.}
Un láser de helio-neón consiste en dos espejos planos y paralelos separados por $L = 35.124\text{ cm}$. El medio gaseoso amplifica luz en un rango espectral de longitudes de onda entre $632.80840\text{ nm}$ y $632.80980\text{ nm}$.
\begin{enumerate}
    \item Encuentre el número de modos longitudinales que se forman en la cavidad y determine la longitud de onda exacta de cada uno (con 8 dígitos significativos).
    \item Calcule la rapidez media cuadrática de un átomo de neón a una temperatura de $120^\circ\text{C}$.
    \item Demuestre que el ancho de banda real debido al ensanchamiento Doppler térmico es consistente con que la cavidad soporte los modos calculados.
\end{enumerate}